% Figure: operation count of Horner's nested evaluation against the naive
% term-by-term evaluation, as a function of polynomial degree. Horner is
% linear in the degree; the naive form that recomputes each power is
% quadratic. The gap is the accuracy-and-speed argument for the nested form.
\begin{figure}[htbp]
\centering
\begin{tikzpicture}
\begin{axis}[
  width=11cm, height=6.5cm,
  xlabel={polynomial degree $n$},
  ylabel={multiplications},
  xmin=0, xmax=20,
  ymin=0, ymax=210,
  grid=both,
  grid style={engblack!12},
  tick label style={font=\footnotesize},
  label style={font=\small},
  legend style={font=\footnotesize, at={(0.03,0.97)}, anchor=north west},
]
% Horner: n multiplications.
\addplot[engblue, very thick] coordinates {
  (0,0) (2,2) (4,4) (6,6) (8,8) (10,10) (12,12) (14,14) (16,16) (18,18) (20,20)
};
% Naive term-by-term, recomputing each power: n(n+1)/2 multiplications.
\addplot[engred, very thick] coordinates {
  (0,0) (2,3) (4,10) (6,21) (8,36) (10,55) (12,78) (14,105) (16,136)
  (18,171) (20,210)
};
\node[font=\footnotesize, engblue, anchor=west] at (axis cs:15,26)
  {Horner: $n$};
\node[font=\footnotesize, engred, anchor=east] at (axis cs:16,150)
  {naive: $\tfrac{n(n+1)}{2}$};
\legend{{Horner (nested)}, {naive (term by term)}}
\end{axis}
\end{tikzpicture}
\caption[Operation count of Horner versus naive evaluation]{Multiplication
count for evaluating a degree-$n$ polynomial. Horner's nested form uses
$n$ multiplications, linear in the degree. The naive term-by-term form
that recomputes each power $x^{k}$ from scratch uses $n(n+1)/2$
multiplications, quadratic in the degree. At degree $20$ the naive form
does $210$ multiplications against Horner's $20$, a factor of about ten.
The speed gap widens with degree; the accuracy gap, from never forming
the large intermediate powers that the naive sum then partly cancels,
is the reason the nested form is the working default.}
\label{fig:vol02:ch01:horner-ops}
\end{figure}
